\documentclass{article}

% Core math + hyperlinks (does not change fonts/margins)
\usepackage{amsmath, amssymb}
\usepackage{hyperref}
\usepackage{graphicx}
\usepackage{booktabs}
\usepackage{todonotes}
\usepackage{float}

% ---- Metadata (edit) ----
\title{CMDO 2025 - Sports Tournament Scheduling
(STS) problem}
\author{%
Luca Anzaldi \texttt{luca.anzaldi@studio.unibo.it} \and
Martino Michelotti \texttt{martino.michelotti@studio.unibo.it} \and
Alex-Cristian Cozma \texttt{alexcristian.cozma@studio.unibo.it}
}
\date{\today}

% ---- Helpful macros (optional) ----
\newcommand{\N}{\mathbb{N}}
\newcommand{\Z}{\mathbb{Z}}
\newcommand{\R}{\mathbb{R}}
\newcommand{\opt}{\mathrm{opt}}
\newcommand{\AllDiff}{\texttt{allDifferent}}

\begin{document}
\maketitle

\section{Introduction}

This report addresses the solution of the well-known \emph{Sports Tournament Scheduling Problem (STS\cite{csplib_prob026})}.
The focus is on presenting different modeling approaches and optimization techniques to handle the scheduling constraints and objectives.  
To avoid repetition, later sections will refer back to this introduction whenever these shared components are required.

\paragraph{Instance parameters}
Let \(T=n\) be the number of teams (even), \(W=T-1\) the number of weeks, and \(P=T/2\) the number of periods (parallel timeslots) per week.
Each solution is modeled referring to these three main variables.

\begin{itemize}
    \item Example1 : $ T = 6 \xrightarrow{} W = 4 $ and  $ P = 3 $
    \item Example2 : $ T = 16 \xrightarrow{} W = 15 $ and  $ P = 8 $
    \item Example3 : $ T = 5  \xrightarrow{} $ NOT VALID 
\end{itemize}
For the sake of brevity, whenever $w$, $p$ and $t$ (or eventually $t_1, t_2$) are used in formulas and the domains are omitted, they have to be implicit meant to be: 
$w \in \{1,\dots, W\}$, $p \in \{1,\dots, P\}$ and $t \in \{1,\dots, T\}$.

\paragraph{General constraints}
The classical formulation of the STS problem requires that:
\begin{enumerate}
    \item every team plays with every other team exactly once;
    \item every team plays exactly once per week;
    \item every team plays at most twice in the same period over the whole tournament.
\end{enumerate}

\paragraph{Objective function}

The main optimization criterion is to achieve a balanced schedule,
in which every team plays approximately the same number of home and away
games. 
The objective is to minimize the overall imbalance:

\[
\sum_{t=1}^{T} \Big| \; 
\#\text{home\_games}(t) - \#\text{away\_games}(t) \; \Big|
\]

Since $W$ is uneven by the definition of the problem, it is clear that no team could ever play the same number of home and away games and therefore the minimal value of $\big|  
\#\text{home\_games}(t) - \#\text{away\_games}(t)  \big|$ for every team $t$ cannot be lower than 1.
This implies that the objective function is (non-strictly) bounded below by the number of teams $T=n$, that we will find out that it corresponds to the optimal value we are looking for.



% --------------------------------- SAT ------------------------------------------------------
\section{SAT Model}

The SAT solution was made using two different modeling approaches. Both are included in the report. 

\vspace{1em}
\textit{Model~1} was implemented using several alternative SAT encodings for the cardinality
constraints introduced above. In particular, we considered the \textbf{Bitwise} encoding,
the \textbf{Heule} encoding, the \textbf{Pairwise} encoding, and the \textbf{Sequential} encoding.
Each encoding provides a different trade-off between the number of variables,
the number of clauses, and the strength of constraint propagation.

\vspace{1em}
\textit{Model~2} extends the formulation of Model~1 by adding further constraints and auxiliary decision variables. In contrast to Model~1, the implementation adopts a single SAT encoding approach (e.g., pairwise-style cardinality constraints) and does not perform a systematic comparison across alternative encodings.  

Moreover, the implementation can optionally enable a \textbf{precomputing} phase described in paragraph \ref{par:precomputing_sat}.





\subsection{Decision variables}

\subsubsection{Model - 1}

The use of multiple encodings within the same model allows for a systematic
comparison of their impact on solver performance, while keeping the underlying
problem formulation fixed.

\paragraph{Booleans}

\begin{itemize}
    \item  $ X_{t,h,p,w}\in\{0,1\} \quad\text{meaning team }t\text{ is scheduled with role }h\text{ in slot }(p,w). $
\end{itemize}

\subsubsection{Model - 2}

\paragraph{Booleans}

\begin{itemize}
    \item  $ M_{t_1,t_2,w} \in \{0,1\} \quad\text{meaning teams $t_1$ and $t_2$ play against each other in week $w$}. $
    \item $HOME_{t,w} \in \{0,1\} \quad\text{meaning team $t$ plays at home in week $w$ (away otherwise)}.$
    \item $P_{t,p,w} \in \{0,1\} \quad\text{meaning team $t$ is assigned to period $p$ in week $w$}.$
\end{itemize}


\paragraph{Remarks.}

Model--2 introduces explicit pair variables \(M_{t_1,t_2,w}\) to capture
which two teams meet in a given week, home assignment variables
\(HOME_{t,w}\) to distinguish roles, and period variables
\(P_{t,p,w}\) to ensure consistency of scheduling.


\subsection{Objective function}

\subsubsection{Model - 1}

We minimize total home/away imbalance across teams:
\[
%\min
\;\; \sum_{t=1}^{T} \left| 
\sum_{p=1}^{P}\sum_{w=1}^{W} \big( X_{t,0,p,w} - X_{t,1,p,w} \big)
\right|.
\]
In the implementation with \texttt{Z3}, the optimization is performed by 
\emph{iterative (binary) search} on a pseudo-Boolean bound \(z\). 
At each step we introduce a constraint of the form 
\(\sum_t \lvert \cdot \rvert \leq z\), encoded using cardinality 
constraints, and solve for feasibility. If satisfiable, the bound 
is tightened; otherwise the process stops at the last feasible value.

\subsubsection{Model - 2 }

As in Model--1, the optimization criterion is to balance the number of
home and away games for each team. Formally:

\[
%\min 
\sum_{t=1}^T \left| \sum_{w=1}^W HOME_{t,w}
      - \Big(W - \sum_{w=1}^W HOME_{t,w}\Big) \right|.
\]

% ======= CONSTRAINT =======
\subsection{Constraints}

\subsubsection{Model - 1 }

\paragraph{Main constraints.}

\begin{itemize}
\item[(C1)] \textbf{Every pair of teams plays at most once:}  
%For each distinct pair of teams $t_1,t_2 \in \{1,\dots,T\}$:
\[
\forall\, t_1,t_2 \;\;\text{with}\;\;t_1\neq t_2\quad
\sum_{p=1}^{P}\sum_{w=1}^{W}
\Big( \bigvee_{h\in\{0,1\}} X_{t_1,h,p,w} \;\land\;
       \bigvee_{h\in\{0,1\}} X_{t_2,h,p,w} \Big)
\;\;\leq 1.
\]

\item[(C2)] \textbf{Each team plays exactly once per week:}  
%For every team $t$ and week $w$:
\[
\forall \,t, w\;
\sum_{h\in\{0,1\}}\sum_{p=1}^{P} X_{t,h,p,w} \;=\; 1.
\]

\item[(C3)] \textbf{Each team appears at most twice in the same period over the tournament:}  
%For every team $t$ and period $p$:
\[
\forall\, t, p\;
\sum_{h\in\{0,1\}}\sum_{w=1}^{W} X_{t,h,p,w} \;\leq\; 2.
\]
\end{itemize}

\paragraph{Implied constraints.}

\begin{itemize}
\item[(C4)] \textbf{Each game slot has exactly one home and one away team:}  
%For every period $p$ and week $w$:
\[
\forall\, p, w\quad
\sum_{t=1}^{T} X_{t,1,p,w} = 1
\quad\text{and}\quad
\sum_{t=1}^{T} X_{t,0,p,w} = 1.
\]
\end{itemize}

\subsubsection{Model - 2 }

\paragraph{Main constraints.}

\begin{itemize}
\item[(C1)] \textbf{Each team plays with every other team only once;}
\[
\forall\, t_1<t_2\qquad
\sum_{w=1}^{W} M_{t_1,t_2,w} \;=\; 1.
\]

\item[(C2)] \textbf{Each team plays exactly once per week}
\[
\forall\, t,\,w\qquad
\sum_{p=1}^{P} P_{t,p,w} \;=\; 1.
\]

\item[(C3)] \textbf{Each team plays at most twice in the same period}
\[
\forall\, t,\,p\qquad
\sum_{w=1}^{W} P_{t,p,w} \;\le\; 2.
\]

\item[(C4)] \textbf{Each slot \((p,w)\) hosts exactly two teams.}
\[
\forall\, p,\,w\qquad
\sum_{t=1}^{T} P_{t,p,w} \;=\; 2.
\]
\end{itemize}

\paragraph{Implied constraints.}

\begin{itemize}
\item[(C5)] \textbf{Home/away consistency when two teams meet in week \(w\).}
\[
\forall\, w,\, t_1<t_2\qquad
M_{t_1,t_2,w} \;\Rightarrow\; \big(HOME_{t_1,w} \oplus HOME_{t_2,w}\big).
\]

\item[(C6)] \textbf{Link between pair variables and period assignments.}
\[
\forall\, w,\, t_1<t_2\qquad
M_{t_1,t_2,w} \;\iff\; \bigvee_{p=1}^{P}\big(P_{t_1,p,w}\land P_{t_2,p,w}\big),
\]
%\[\forall\, w,\,p,\, t_1<t_2\qquad\big(P_{t_1,p,w}\land P_{t_2,p,w}\big)\Rightarrow\; M_{t_1,t_2,w}.\]
%(Equivalently: \(M_{t_1,t_2,w}\) iff the two teams share some period \(p\) in week \(w\).)
\end{itemize}

\paragraph{Precomputing tecnique.}
\label{par:precomputing_sat}

\begin{itemize}
\item[(P1)] \textbf{Precomputed round-robin pruning.}
Let \(\mathcal{R}_w\subseteq\{\{t_1,t_2\}\mid t_1<t_2\}\) be the fixed set of pairs for week \(w\) produced by the round-robin generator. Then:
\[
\forall\, w,\, t_1<t_2\qquad
M_{t_1,t_2,w} =
\begin{cases}
1, & \text{if }\{t_1,t_2\}\in \mathcal{R}_w,\\[4pt]
0, & \text{otherwise.}
\end{cases}
\]
\end{itemize}


% ======= VALIDATION SAT =======

\subsection{Validation}
\label{sub:validation_sat}

\paragraph{Experimental Setup} \mbox{}\\
\\
All experiments were conducted using the \texttt{Z3} solver for SAT and
optimization (\cite{z3}). The implementation is provided in Python, and the solver is
invoked through the \textbf{Z3 API}. 
\\
%The experiments were tested on \emph{Linux} on a laptop equipped with a
%\textit{12th Gen Intel(R) Core(TM) i7-1280P} processor and an
%\textit{NVIDIA RTX-3060 GPU}. No GPU acceleration was employed for the SAT
%solving itself, which ran entirely on the CPU.

Reproducibility instructions, including installation steps and execution
commands, are provided in the accompanying \texttt{README.md} file.



\paragraph{Experimental results} \mbox{}\\
\label{par:Experimental_result_SAT}

The \textbf{following tables reports} the results of the SAT experiments. 
Each row corresponds to a different value of $n$. 
On the tables are reported tthe value of the objective value obtained using different approach and in brackets the computation time for the best solution found. If the instance is solved to optimality is in bold. In addition if the instance is proved to be
unsatisfiable, it is indicate as \texttt{UNSAT} and if no answer is obtained within the
time limit of 300s are indicate it with \texttt{N/A}.
Each modeling approach have is own table (Tab.~\ref{tab:results_sat_m1} and Tab.~\ref{tab:results_sat_m2}).  

The \textbf{precomputation time} required to generate the round-robin pairs is not reported in the tables, as it is negligible: it remains below one second even for 
$n\approx50$, and therefore has no practical impact on the overall computational performance.

% MODEL 1 and MODEL 2 - UPDATED TABLE
\begin{table}[h]
\centering
\scriptsize

\begin{minipage}[t]{0.55\textwidth}
\centering
\setlength{\tabcolsep}{3pt}
\renewcommand{\arraystretch}{1.05}
\begin{tabular}{c|c|c|c|c}
\hline
$n$ & Bitwise & Heule & Pairwise & Sequential \\
\hline
2   & \textbf{2} (0s)  & \textbf{2} (0s)  & \textbf{2} (0s)  & \textbf{2} (0s) \\
4   & \texttt{UNSAT} (0s) & \texttt{UNSAT} (0s) & \texttt{UNSAT} (0s) & \texttt{UNSAT} (0s) \\
6   & \textbf{6} (0s)  & \textbf{6} (0s)  & \textbf{6} (0s)  & \textbf{6} (0s) \\
8   & 14 (1s)          & 24 (1s)          & 18 (1s)          & 24 (0s) \\
10  & 26 (6s)          & 36 (4s)          & 18 (7s)          & 20 (6s) \\
12  & 30 (87s)         & 44 (13s)         & 36 (40s)         & 30 (151s) \\
14  & \texttt{N/A}     & \texttt{N/A}     & \texttt{N/A}     & \texttt{N/A} \\
16  & \texttt{N/A}     & \texttt{N/A}     & \texttt{N/A}     & \texttt{N/A} \\
\hline
\end{tabular}
\caption{Experimental results of Model 1.}
\label{tab:results_sat_m1}
\end{minipage}
\hfill
\begin{minipage}[t]{0.42\textwidth}
\centering
\setlength{\tabcolsep}{3pt}
\renewcommand{\arraystretch}{1.05}
\begin{tabular}{c|c|c}
\hline
$n$ & Standard & Precomputing \\
\hline
2   & \textbf{2} (--)      & \textbf{2} (0s) \\
4   & \texttt{UNSAT} (--)  & \texttt{UNSAT} (0s) \\
6   & \textbf{6} (--)      & \textbf{6} (0s) \\
8   & \textbf{8} (--)      & \textbf{8} (0s) \\
10  & \textbf{10} (--)     & \textbf{10} (0s) \\
12  & \textbf{12} (3s)     & \textbf{12} (1s) \\
14  & \textbf{14} (83s)    & \textbf{14} (6s) \\
16  & 46 (244s)            & \textbf{16} (10s) \\
18  & \texttt{N/A}         & 44 (53s) \\
20  & \texttt{N/A}         & 200 (272s) \\
22  & \texttt{N/A}         & \texttt{N/A} \\
\hline
\end{tabular}
\caption{Experimental results of Model 2.}
\label{tab:results_sat_m2}
\end{minipage}
\end{table}



% --------------------------------- MIP ------------------------------------------------------
\section{MIP Model}

In this section we will discuss how we implemented a Mixed Integer Linear Program model able to solve the given Tournament Scheduling problem.

The main obstacle in the implementation of linear programs is the fact that the objective function and every constraint have to be linear in our variables; therefore we transformed logic propositions and non-linear functions with linear inequalities, introducing, if needed, some auxiliary variable as for the implementation of the objective function.

\subsection{Decision variables}
Before defining the variables, we compute the list $L$ of the possible match-ups, without caring about the order of the $(i,j)$ team tuples (we impose $i<j$ for unicity).
If $n=T$ is the number of teams, we have $|L|=\binom{n}{2} = \frac{n!}{2!(n-2)!}$, i.e. $\binom{n}{2}$ games to be played.


Our decision variables are then two 3-dimensional tensors of \textbf{boolean} values $Y_{w,p,(i,j)}$ and $H_{w,p,(i,j)}$, where
\begin{itemize}
    \item $Y_{w,p,(i,j)} = 1 \iff$ teams $i$ and $j$ are playing one against the other in week $w$ in period $p$;
    \item $H_{w,p,(i,j)} = 1 \iff$ team $i$ (the team with lowest index since $\forall (i,j) \in L$, $i<j$) is playing home in week $w$ in period $p$.
\end{itemize}

Clearly these variables have to be linked one with the other, and that's why we will introduce some variable-linking constraint to our model.


\subsection{Objective function}

The objective function is used only in the optimization version of the MIP
model. Its purpose is to balance the number of home and away games played by
each team.

For each team $t$, let $h_t$ be the number of home games played by team $t$.
Since each team plays exactly one match per week and there are $W = n - 1$
weeks, the number of away games is:

\[
a_t = W - h_t.
\]

The home/away imbalance of team $t$ is therefore:

\[
|h_t - a_t|.
\]

Since absolute values cannot be used directly in a linear objective function, we
introduce an auxiliary integer variable $b_t$ for each team $t$, representing
the imbalance of that team. We impose the following linear constraints:

\[
b_t \geq h_t - a_t \qquad \forall t,
\]

\[
b_t \geq a_t - h_t \qquad \forall t.
\]

Since the objective minimizes the sum of the imbalance variables, each $b_t$
will take the smallest value satisfying both inequalities, and therefore:

\[
b_t = |h_t - a_t|.
\]

The objective function of the optimization version is then:

\[
\min \sum_{t=1}^{n} b_t.
\]

In the decision version, no real objective function is optimized. A dummy
constant objective is used only to run the feasibility model through the MIP
solver.



% ======= CONSTRAINTS MIP =======
\subsection{Constraints}

\paragraph{Variable linking constraint}
As already mentioned, we have to link the variables and this has to be done by adding linear inequalities to our linear program.

The only constraint we have to add is the one that ensures that $H_{w,p,(i,j)}\neq1$ if teams $i$ and $j$ do not play against in week $w$ in period $p$, i.e.
\[
Y_{w,p,(i,j)} = 0\implies H_{w,p,(i,j)}\neq1.
\]
This constraint can be translated in a linear inequality by imposing
\[
H_{w,p,(i,j)} \leq Y_{w,p,(i,j)}.
\]

\paragraph{Main constraints.}


\textbf{(C1) Each pair of teams plays exactly once:}

\[
\forall (i,j) \in L \qquad
\sum_{w=1}^{W} \sum_{p=1}^{P} Y_{w,p,(i,j)} = 1.
\]

\textbf{(C2) Each team plays exactly once per week:}

\[
\forall t,w \qquad
\sum_{p=1}^{P}
\sum_{\substack{(i,j)\in L \\ i=t \vee j=t}}
Y_{w,p,(i,j)} = 1.
\]

\textbf{(C3) Each team plays at most twice in the same period:}

\[
\forall t,p \qquad
\sum_{w=1}^{W}
\sum_{\substack{(i,j)\in L \\ i=t \vee j=t}}
Y_{w,p,(i,j)} \leq 2.
\]

\textbf{(C4) Each slot $(w,p)$ hosts exactly one game:}

\[
\forall w,p \qquad
\sum_{(i,j)\in L} Y_{w,p,(i,j)} = 1.
\]

With this model, we do not need to explicitly impose that exactly two teams
play in each slot, since the possible games have already been represented as
team pairs in $L$.

\paragraph{Symmetry breaking constraints}
\begin{itemize}
    
    \item[(C5)] We fix each match in the first week by imposing:
    \begin{align*}
        \forall p =1,\dots,P\qquad   &Y_{1,p,(2p-1,\,2p)}=1\\
        &H_{1,p,(2p-1,\,2p)}=1
    \end{align*}
    
    \item[(C6)] Since the weeks are interchangeable in our formulation (constraints and objective do not depend on the week index), we can break the week-permutation symmetry by fixing the opponent of team 1 in each week, without fixing the period.
    \[
    \forall w = 1,\dots,W \qquad \sum_{p = 1}^P Y_{w,p,(1,\,w+1)} = 1,
    \]
    i.e., in week \(w\) team \(1\) plays against team \(w+1\) in some period \(p\).

    \item[(C7)] In addition we break symmetry by fixing the main diagonal of the schedule matrix: in each diagonal slot \((w,p)=(p,p)\) (for the first \(P\) weeks), team \(1\) must be scheduled, without fixing the opponent. Formally:
    \[
    \forall p = 1,\dots,P \qquad
    \sum_{j=2}^{T} Y_{p,p,(1,j)} = 1.
    \]
    
\end{itemize}

% ======= VALIDATION MIP =======
\subsection{Validation}

\paragraph{Experimental setup.}
All experiments for the MIP formulation were conducted using the CBC solver
through the \texttt{python-mip} library~\cite{pythonmip}. The implementation is
provided in Python, and the solver is invoked through the \texttt{python-mip}
API.

For each instance, two versions of the model were executed:
\begin{itemize}
    \item \texttt{mip\_cbc\_decision}, corresponding to the decision version of
    the problem, where the goal is to find any feasible schedule;
    \item \texttt{mip\_cbc\_optimization}, corresponding to the optimization
    version, where the objective is to minimize the total home/away imbalance.
\end{itemize}

The runtime is measured as total wall-clock time, including both model
construction and solving time. A global time limit of 300 seconds is imposed.
For solved instances, the reported time is the floor of the actual total runtime.
If the decision version does not find a feasible schedule within the time limit,
the runtime is reported as N/A. Similarly, for the optimization version, the
runtime is reported only when optimality is proved within the time limit.

For optimization runs that do not prove optimality within the time limit, the
objective value, when available, reports the best solution quality found by the
solver before timeout.

\paragraph{Experimental results.}

\begin{table}[H]
\centering
\begin{tabular}{c|cc|ccc}
\hline
& \multicolumn{2}{c|}{Decision}
& \multicolumn{3}{c}{Optimization} \\
\cline{2-6}
$n$ & Solved & Time (s)
& Optimal & Time (s) & Obj. \\
\hline
2  & Yes & 0   & Yes & 0   & 2  \\
4  & UNSAT & N/A & UNSAT & N/A & N/A \\
6  & Yes & 0   & Yes & 0   & 6  \\
8  & Yes & 0   & Yes & 0   & 8  \\
10 & Yes & 6   & Yes & 13  & 10 \\
12 & Yes & 87  & Yes & 270 & 12 \\
\hline
\end{tabular}
\caption{Experimental results of the MIP model using CBC. For decision runs,
``Solved'' indicates whether a feasible schedule was found within 300 seconds.
For optimization runs, ``Optimal'' indicates whether optimality was proved within
300 seconds. The runtime is reported as N/A when no feasible or optimal solution
is available for the corresponding setting. Objective values refer to the
home/away imbalance minimized by the optimization model.}
\label{tab:mip_results}
\end{table}



% ---------------- CP ----------------
\clearpage
\section{CP Model}

In this last section we will discuss how we implemented the Constraint programming model.

With respect to the other two implementations, constraint programming admits a higher range of modeling possibilities and we tried to use them to model less sparse arrays, by avoiding the heavy use of boolean variables as we did in the other models.

\subsection{Decision variables}
Our decision variables are 3 different arrays:
\begin{itemize}
    \item $P_{t,w} \in \{1, \dots,P\}$ where 
    \begin{center}
    $ P_{t,w} = p \iff$ team $t$ plays in period $p$ in week $w$
    \end{center}
    \item Boolean array $H_{t,w}$ where 
    \begin{center}
    $ H_{t,w} = \verb|true| \iff$ team $t$ plays home in week $w$;
    \end{center}
    \item $OPP_{t_1,w}\in \{1,\dots,T\}$ where
    \begin{center}
    $OPP_{t_1,w} = t_2 \iff$ team $t_1$ plays against team $t_2$ in week $w$.
    \end{center}
\end{itemize}

\subsection{Objective function}\label{CP obj function}
For the objective function, we opted for a proxy of the one defined in the introduction.
As discussed in section \ref{obj function MIP}, we used the \verb|max| function instead of \verb|abs| and the Home/Away imbalance was parametrized using just home games; but in this case, instead of summing up the imbalance for each team, we decided to look for the maximum value of imbalance among all teams and this turned out to be a good choice for performance improvement.
The objective function in this section will be
\[
\max_{t \in \{1,\dots, n\}} \big| \; 
\#\text{home\_games}(t) - \#\text{away\_games}(t) \; \big|.
\]
Such a function can be used instead of the original one since we are looking to the optimal value: $n$ for the original formulation, 1 for the $max$ objective function.
Knowing a hypothetic value $k_{max}$ of the $max$ objective function, defines an upper bound on the values $k_{OG}$ of the original objective function, in fact
\[
k_{OG}\,\leq\, n\,k_{max} \qquad\forall\, n, k_{max}
\]
and in case of optimality we clearly get
\[
k_{max}=1 \implies k_{OG} \leq n \implies k_{OG}=n.
\]


\subsection{Constraints}


\paragraph{Main constraints.}
\begin{itemize}
\item[(C1)] \textbf{Every pair of teams plays at most once.}  
For this constraint we used the global constraint \verb|alldifferent| in the following way:
\[
\forall\, t \qquad \verb|alldifferent|(\,\{\,OPP_{t,w}\,|\, w = 1,\dots, n-1\,\}\,).
\]

\item[(C2)] \textbf{Each team plays exactly once per week.}  

This constraint has not to be defined in this model because the array of variables $P_{t,w}$ admits just one value for every team in every week.

\item[(C3)] \textbf{Each team appears at most twice in the same period over the tournament.}  
\[
\forall t, p \qquad \big|\{w\,| P_{t,w}=p\}\big| \leq 2
\]
\end{itemize}


\paragraph{Implied constraints.}
These constraints are crucial for the actual definition of the problem and the correct use of the decision variables
\begin{itemize}
\item[(C4)] \textbf{No team plays against itself:}
\[
\forall t , w \qquad OPP_{t,w} \neq t
\]

\item[(C5)] \textbf{Matches are symmetric:}
\[
\forall t , w \qquad OPP_{OPP_{t,w}\,, w} = t
\]

\item[(C6)] Since the weeks are interchangeable in our formulation (constraints and objective do not depend on the week index), we can break the week-permutation symmetry by fixing the opponent of team 1 in each week, without fixing the period.
\[
\forall w = 1,\dots,W \qquad \sum_{p = 1}^P Y_{w,p,(1,\,w+1)} = 1,
\]
i.e., in week \(w\) team \(1\) plays against team \(w+1\) in some period \(p\).

\item[(C7)] In addition we break symmetry by fixing the main diagonal of the schedule matrix: in each diagonal slot \((w,p)=(p,p)\) (for the first \(P\) weeks), team \(1\) must be scheduled, without fixing the opponent. Formally:
\[
\forall p = 1,\dots,P \qquad
\sum_{j=2}^{T} Y_{p,p,(1,j)} = 1.
\]

\end{itemize}


\paragraph{Symmetry breaking}
The problem has different symmetries and we tried to reduce the search space as much as possible, but maintaining some flexibility. Some of the subsequent constraints seem redundant, but we experimentally found out they were improving performance by helping propagation.
The constraints are imposed especially on team 1 and week 1 for simplicity, but they could've been imposed on every other week and team (singularly) because of the week and team symmetries.
\begin{itemize}
    \item [(C9)] We fix each match in the first week by imposing three combined constraints:
    \begin{itemize}
        \item fixing match-ups:
            \[
            \forall p \qquad OPP_{2p-1, 1} = 2p \; \land \;OPP_{2p, 1} = 2p-1\;;
            \]
        \item fixing periods:
            \[
            \forall p \qquad P_{2p-1, 1} =  p \; \land \; P_{2p, 1}\;;\qquad\qquad\qquad
            \]
        \item fixing home/away teams:
            \[
            \forall p \qquad H_{2p-1, 1} =  \verb|true| \; \land \; H_{2p, 1} = \verb|false|\;.
            \]
    \end{itemize}

    \item [(C10)] We impose an increasing order of the opponents for team 1 with the following constraint:
    \[
    \forall w = 1,\dots, n-2\qquad OPP_{1,w} < OPP_{1, w+1}.
    \]

    This constraint helps, along with (C9), to break the team symmetry of the problem by defining a total order on the teams' labels.
    
\end{itemize}


\subsection{Validation}



\paragraph{Experimental design.}
All experiments for the constraint programming model were done using MiniZinc IDE and then implemented on the docker container. The model is runnable using Gecode solver, but with poor performance; we had some slight improvement when used with Optimization level O2 (two pass compilation).
The model was instead optimized for the Chuffed solver (v 0.13.2), since we noticed it's efficiency already from the early stages of modeling.
For reproducibility we fixed the random seed 42.

The experiment were ran on a \textit{Windows 11} OS and the hardware platform was a laptop equipped with a \textit{AMD RYZEN 7 5800H} processor. No GPU acceleration was employed for CP solvers, which ran entirely on
the CPU.  \\



\paragraph{Experimental results.}
The interpretation of the reported results and the meaning of the table entries follow a similar convention to the one described in Section~\ref{par:Experimental_result_SAT}, with the difference that we reported here the proxy objective function described in Section ~\ref{CP obj function}.
\begin{table}[h]
\centering
\caption{Results with/without symmetry breaking}
\label{tab:cp}
\begin{tabular}{lcccc}
\toprule
\textbf{$n$} & \textbf{Gecode+SB} & \textbf{Gecode--SB} & \textbf{Chuffed+SB} & \textbf{Chuffed--SB}\\
\midrule
2  & \textbf{1} & \textbf{1} & \textbf{1} & \textbf{1} \\
4  &\texttt{UNSAT}&\texttt{UNSAT}&\texttt{UNSAT}&\texttt{UNSAT}\\
6  & \textbf{1}  &\textbf{1}  & \textbf{1} & \textbf{1} \\
8  & \textbf{1}  & \textbf{1} & \textbf{1} & \textbf{1} \\
10 & \textbf{1}  &\texttt{N/A}& \textbf{1} & \textbf{1} \\
12 & \texttt{N/A}&\texttt{N/A}& \textbf{1} & \textbf{1} \\
14 & \texttt{N/A}&\texttt{N/A}& \textbf{1} & \textbf{1} \\
16 & \texttt{N/A}&\texttt{N/A}& \textbf{1} & \textbf{1} \\
18 & \texttt{N/A}&\texttt{N/A}& 13         & \textbf{1}     \\
20 & \texttt{N/A}&\texttt{N/A}& 19         &\texttt{N/A}\\
22 & \texttt{N/A}&\texttt{N/A}& \texttt{N/A} &\texttt{N/A}\\
\bottomrule
\end{tabular}
\end{table}




% --------------- Conclusions ---------------
\section{Conclusions}
The diversity of the programming languages and their limitations and strengths gave us different views on the problem and tools to handle it, letting us develop different models, having diverse variables and therefore also diverse constraints; without limiting ourselves to just develop three implementations of the same model. 

Among all models we developed, unless one specific case ($n$=18 with Constraint Programming), we had a good improvement for each of them when symmetry breaking (SB) constraints were applied.
For all three models, the respective SB constraints allowed us to find a solution when the problem was scaled further by increasing $n$ by one (+2) step.
Execution times were not reported, but they also had an improvement by adding SB constraints, outputting faster solutions by reducing the search tree.

We initially implemented constraints that fixed the home and away balances to directly force them on the optimal values. The mentioned constraints are
\[
\sum_{t=1}^{n}\big| \; 
\#\text{home\_games}(t) - \#\text{away\_games}(t) \; \big| = n.
\]
or alternatively
\[
\forall\, t \qquad \big| \; 
\#\text{home\_games}(t) - \#\text{away\_games}(t) \; \big| = 1
\]

By strictly constraining the objective function, the model was finding a solution (optimal by problem definition) in significant less time; but this would've changed the optimization problem in a satisfaction problem, so we kept them away to let the solver find the optimal solution without external hints.

The \textbf{SAT} approach has a clear limitation in its restriction to purely Boolean variables, which significantly constrained how we could define decision variables. Despite the considerable effort of implementing and testing four different cardinality encodings, the improvements in the final results remained limited for our formulation.
The SAT experiments highlighted that, in our case, performance is influenced more by the modeling choices than by symmetry-breaking techniques. While symmetry breaking yielded incremental improvements, reformulating the model produced substantially larger performance gains.

\textbf{Linear Programming} allowed us to set integer variables in the problem, but it would have been trickier to define the constraints, especially when summing up variables. Further issues encountered with MIP were the linear limitations we had defining functions and constraint, bringing us to define new auxiliary variables and reformulate them as linear inequalities.

\textbf{Constraint programming} was for sure the more flexible method, admitting in the formulation every function we needed and different kind of variables; but it's exactly here where we had to think and experiment plenty of different options to find the best combination that could facilitate the constraint propagation of the solver and consequently the reduction of the search tree.

% --------------- Authorship ---------------
\newpage
\section*{Authenticity and Author Contribution Statement}
We declare that the work reported here is our own and properly cites all external ideas and sources.
Parts of this project (text/code/experiments) were assisted by AI tools where indicated; we disclose the tools used and the exact sections where AI assistance was applied expect for the part where it was used to improve the quality of the report, such as grammar, sentence formulation and for code debugging.
\\[0.5em]

\paragraph{Author contributions.}
Luca Anzaldi developed the SAT-based formulations, including both Model~1 and
Model~2. He was responsible the
configuration and setup of all scripts, Docker environments, and auxiliary
tools required to run the experiments.
Martino Michelotti developed the complete CP model.
The report was written by both students, mostly reporting the part of the project they were responsible of.

% --------------- References ---------------
\begin{thebibliography}{9}

\bibitem{minizinc}
MiniZinc: The Modeling Language.\\
\url{https://www.minizinc.org/}

\bibitem{csplib_prob026}
CSPLib, \emph{Problem 026: Sports Tournament Scheduling (prob026)}.\\
\url{https://www.csplib.org/Problems/prob026/}
(Accessed: 2026-01-29)

\bibitem{z3}
Leonardo de Moura and Nikolaj Bj{\o}rner.
\emph{Z3: An Efficient SMT Solver}.
In \emph{Proceedings of the 14th International Conference on Tools and Algorithms for the Construction and Analysis of Systems (TACAS)}, 
LNCS 4963, pp. 337--340, Springer, 2008.\\
\url{https://github.com/Z3Prover/z3}

\bibitem{pythonmip}
Henrique R. Lopes, Thiago A. F. S. Santos, and Tallys H. Y. Lopes.
\emph{Python-MIP: A Python Toolbox for Modeling and Solving Mixed-Integer Linear Programs}.
SoftwareX, Volume 12, 2020.\\
\url{https://www.python-mip.com/}


\end{thebibliography}

\end{document}
